% Propietario conservado: /Users/ruben/Documents/New project/output/EXTREMA_MEDIA_RAZON_RECIPROCIDAD_ES_EN_20260918/ARTICULO/ES/source/sections/generacion.tex
\section{Generación coinductiva de \texorpdfstring{\(\pi,\varphi,e\)}{pi, phi, e} a profundidad arbitraria}
\label{x_reciprocidad:sec:generacion}

La emisión de seis bloques es una etapa finita de la construcción, no una expansión decimal completa. Su función consiste en producir regiones con multiplicidad, orientación y reglas de transporte. La prolongación debe conservar esos datos y determinar una familia compatible de prefijos de longitud arbitraria. La palabra \emph{generación} designa esta operación; \emph{genealogía} designa la reconstrucción de sus dependencias. Ambas nociones intervienen, pero cumplen funciones distintas.

\subsection{Órbitas, regiones y caracteres de clausura, propagación y autoescala}

Sobre una palabra \(w=(w_1,\ldots,w_6)\), definimos
\[
 r(w)=(w_3,w_1,w_2,w_5,w_6,w_4),\qquad
 s(w)=(w_4,w_5,w_6,w_1,w_2,w_3).
\]
Se cumple \(r^3=s^2=I\) y \(srs=r^{-1}\); por tanto estas permutaciones realizan el grupo diedral \(D_3\). La acción conserva el catálogo y las multiplicidades y conmuta con \(\Psi\). El cociente de las \(243\) palabras visibles tiene \(43\) órbitas: treinta y ocho de cardinal seis y cinco de cardinal tres.

La clase de clausura se reconoce por un predicado de fibra: sus seis orientaciones tienen cinco emisiones por palabra, multiplicidad total \(1008\) y distribución \(432+4\cdot144\). Existe una sola órbita con esas propiedades en el catálogo. Para precisar las otras dos selecciones, escribimos \(f(w)=\#\Psi^{-1}(w)\), \(m(w)=\sum_{U\in\Psi^{-1}(w)}\operatorname{mult}(U)\) y \(\nu(w)=(n_0,n_1,n_2)\). La clase diagonal aditiva de una emisión es \(d(U)=[i+j]_9\) en los índices de cero a ocho; su firma aditiva conserva esta clase y la orientación \(ES\) o \(NO\). Sea \(\mathcal D(\mathcal O)\) el conjunto de esas clases para todas las emisiones sobre una órbita.

Entre las órbitas de tamaño seis, el predicado conjunto
\[
 f(w)=1,\qquad m(w)=144\quad(w\in\mathcal O),\qquad
 \mathcal D(\mathcal O)=\{0,3,6\}
\]
deja exactamente seis órbitas. En este sector, el censo \(\nu=(2,3,1)\), igual al de clausura, selecciona únicamente la propagación; el censo \(\nu=(2,2,2)\) selecciona únicamente la autoescala. El soporte diagonal no puede omitirse: sin él hay cuatro órbitas unipuntuales equilibradas de multiplicidad \(144\), no una. Finalmente, la existencia de una emisión con \(d=6\) y orientación \(ES\) selecciona un representante de cada órbita y conduce a
\begin{equation}
 w^{\rm cl}=010211,\qquad w^{\rm pr}=201101,\qquad w^{\rm au}=121200.
 \label{x_reciprocidad:eq:palabras-regionales}
\end{equation}
Estos símbolos designan primero regiones del catálogo. La identificación escalar se demostrará mediante sus lectores; no se eligen porque coincidan con cifras de una constante conocida.

La palabra de clausura \(010211\) tiene las cinco preimágenes siguientes:
\begin{center}\small
\begin{tabular}{@{}lll r@{}}\toprule
Emisión de seis bloques&Firma multiplicativa&Papel&Multiplicidad\\\midrule
\(501,614,498,169,272,272\)&\(555555\)&transversal&432\\
\(810,923,870,169,272,272\)&\(339555\)&orientación positiva&144\\
\(810,983,810,169,272,272\)&\(393555\)&orientación positiva&144\\
\(870,923,810,169,272,272\)&\(933555\)&orientación positiva&144\\
\(870,923,810,418,674,521\)&\(933717\)&retorno orientado&144\\\bottomrule
\end{tabular}
\end{center}
La multiplicidad \(432\) y la firma constante \(555555\) caracterizan la región transversal. No identifican la posición central \((5,5)\) de APP. Confundir una firma de seis ventanas con una posición del soporte eliminaría la distinción entre región numérica y estructura electrónica.

La microfibra posee además una morfología bipartita precisa. Cada emisión es un producto de una clase de cursor aditivo y una clase multiplicativa. Las tres regiones positivas comparten la misma clase aditiva de dieciocho condiciones; juntas contienen tres clases multiplicativas de ocho condiciones cada una. La región transversal posee un producto \(18\times24\); las tres positivas reunidas, otro \(18\times24\); y el retorno, un producto \(18\times8\). La descomposición por emisiones tiene cinco partes, mientras la descomposición por componentes conexas del grafo bipartito tiene tres. Ambas conservan las mismas \(1008\) semillas sin reducir la forma a un cardinal.

% Procedencia: inventario REV05, SR-015--027; ampliacion 04_tpk_emisor,
% firmas y factorizacion; X_RECIPROCIDAD/sections__generacion.tex:16--43.
% Ampliacion de la prueba de metadatos y de las tres componentes.
\subsection{Metadatos orientados y descomposición de la microfibra regional}
\label{sec:microfibra-bipartita-explicita}

La firma aditiva conserva la clase diagonal de origen y la orientación
de su desplazamiento. Para un cursor inicial $(i,j;d)$, con índices entre
cero y ocho, definimos
\[
 c=[i+j]_9,\qquad
 h(d)=\begin{cases}ES,&d\in\{E,S\},\\NO,&d\in\{N,O\}.
 \end{cases}
\]
En las activaciones aditivas consecutivas, la clase $c$ aumenta en uno
para $ES$ y disminuye en uno para $NO$. Tras la novena activación se
invierte la dirección. La sucesión de muestras es, por tanto, determinada
por $(c,h)$, con la inversión posterior a la muestra novena. La muestra
en clase $c'$ es $\rho_9(c'+2)$; cada firma decimal suma tres muestras
consecutivas y reduce el total módulo diez.

\begin{proposition}[Recuperación de los metadatos aditivos]
El lector $(c,h)\mapsto f_+(c,h)$ es biyectivo sobre las dieciocho
firmas aditivas del emisor. Cada preimagen en el espacio de cursores
consta de dieciocho condiciones iniciales. Las dos coordenadas se
recuperan mediante la tabla siguiente.
\end{proposition}

\begin{center}
\begin{tabular}{c|cc}
\toprule
$c$&$f_+(c,ES)$&$f_+(c,NO)$\\
\midrule
0&\texttt{988212}&\texttt{212988}\\
1&\texttt{212645}&\texttt{645212}\\
2&\texttt{546988}&\texttt{988546}\\
3&\texttt{889221}&\texttt{221889}\\
4&\texttt{122564}&\texttt{564122}\\
5&\texttt{465898}&\texttt{898465}\\
6&\texttt{898122}&\texttt{122898}\\
7&\texttt{221456}&\texttt{456221}\\
8&\texttt{654889}&\texttt{889654}\\
\bottomrule
\end{tabular}
\end{center}

\begin{proof}
Para cada entrada se aplica la recurrencia de clases descrita antes del
enunciado durante dieciocho activaciones; sus grupos de tres dan las seis
cifras impresas. Por ejemplo, desde $(c,h)=(4,NO)$ las primeras nueve
clases son $4,3,2,1,0,8,7,6,5$ y las muestras son
$6,5,4,3,2,1,9,8,7$. Los tres totales son $15,6,24$, que publican
$5,6,4$. La inversión deja la segunda mitad $1,2,2$, y resulta
\texttt{564122}. Las dieciocho cadenas de la tabla son distintas,
por lo que el lector es inyectivo y enumera toda su imagen. Para cada
$c$ hay nueve parejas $(i,j)$, una por cada $i$, y para cada $h$ hay
dos direcciones. Esto da $9\cdot2=18$ condiciones por firma y
$18\cdot18=324$ condiciones en total.
\end{proof}

Sean $A_s=f_+^{-1}(s)$ y $B_e=f_\times^{-1}(e)$ las clases de cursores
de las dos hojas. La ley decimal recupera ambas firmas de cada emisión
$U$: si $U_j=100s_j+10t_j+u_j$, entonces
$s_j=\lfloor U_j/100\rfloor$ y $e_j=[t_j-s_j]_{10}$.
En consecuencia, la preimagen de $U$ es exactamente $A_s\times B_e$.
Las condiciones iniciales permanecen en ese producto; sus cardinales
son una lectura adicional de la misma preimagen.

\begin{definition}[Incidencia bipartita de la microfibra]
Para la palabra orientada $w=\texttt{010211}$, se forma el grafo cuyos
vértices de una clase son las firmas aditivas que aparecen sobre $w$ y
cuyos vértices de la otra clase son las firmas multiplicativas. Una
arista $(s,e)$ representa la emisión correspondiente y conserva como
preimagen etiquetada el producto $A_s\times B_e$.
\end{definition}

La recuperación de firmas de las cinco emisiones da la tabla completa
de aristas:
\begin{center}
\begin{tabular}{llrrl}
\toprule
$s$&$e$&$|A_s|$&$|B_e|$&estrato\\
\midrule
\texttt{564122}&\texttt{555555}&18&24&transversal\\
\texttt{898122}&\texttt{339555}&18&8&positivo\\
\texttt{898122}&\texttt{393555}&18&8&positivo\\
\texttt{898122}&\texttt{933555}&18&8&positivo\\
\texttt{898465}&\texttt{933717}&18&8&retorno\\
\bottomrule
\end{tabular}
\end{center}

\begin{theorem}[Descomposición exacta en tres componentes]
La incidencia bipartita tiene tres componentes conexas, isomorfas a
$K_{1,1}$, $K_{1,3}$ y $K_{1,1}$. Su preimagen conjunta es la unión
disjunta
\begin{align*}
 &A_{\texttt{564122}}\times B_{\texttt{555555}},\\
 &A_{\texttt{898122}}\times
  (B_{\texttt{339555}}\sqcup B_{\texttt{393555}}
     \sqcup B_{\texttt{933555}}),\\
 &A_{\texttt{898465}}\times B_{\texttt{933717}}.
\end{align*}
Los cardinales respectivos son $432,432,144$ y suman $1008$.
La descomposición por emisiones consta de cinco productos y la
descomposición por componentes consta de tres conjuntos.
\end{theorem}
\begin{proof}
Las tres firmas aditivas impresas son distintas. Las cinco firmas
multiplicativas también lo son. Sólo las tres aristas positivas comparten
un vértice, el asociado a \texttt{898122}; forman una estrella conexa
con tres aristas. Las otras dos aristas carecen de vértices comunes con
ella o entre sí y cada una forma una componente. Las preimágenes de
firmas diferentes bajo una función son disjuntas. La factorización del
emisor y la distributividad del producto cartesiano sobre esa unión
disjunta prueban las tres expresiones. Sus cardinales son
$18\cdot24$, $18(8+8+8)$ y $18\cdot8$. Las cinco aristas siguen
individualizadas dentro de esas componentes, con su multiplicidad y su
emisión ordenada.
\end{proof}

El recuperador de metadatos da $(c,h)=(4,NO)$ para la transversal,
$(6,ES)$ para cada una de las tres positivas y $(5,NO)$ para el retorno.
Así, orientación, clase diagonal, incidencia y multiplicidad se leen de
las mismas firmas producidas por los cursores. La distinción entre cinco
emisiones y tres componentes conserva la estructura completa del censo,
incluida la coincidencia de cardinal $432$ entre dos componentes distintas.

\subsection{Transporte finito y frontera de prolongación}

Los endomorfismos \(L_0,L_1\) del espacio de palabras ternarias actúan antes de la realización arquimediana. Su determinación comienza por las transiciones del estado aritmético, no por la prescripción de dos matrices. Si \(U_t\) es el paso compuesto de selección, transporte y actualización, el bloque observado satisface
\[
 b_j^\chi=\Pi_6(x_j^\chi),\qquad
 x_{j+1}^\chi=U_{9j+9}\cdots U_{9j+1}(x_j^\chi),
 \qquad \chi\in\{\mathrm{cl},\mathrm{pr},\mathrm{au}\}.
\]
Se forman seis transiciones por cada régimen: dos consecutivas en cada una de las tres regiones. Los pares de matrices de orígenes y destinos obtenidos en el registro son
\[
\begin{array}{c|c|c|c}
X_0&Y_0&X_1&Y_1\\\hline
010211&012222&010211&002111\\
012222&010211&002111&110221\\
201101&121221&102011&012222\\
121221&102011&012222&102011\\
121200&112202&121020&010210\\
112202&121020&010210&010200
\end{array}
\]
Cada cadena de seis dígitos representa una fila en \(\FF_3^6\). Puesto que \(\det X_0=1\) y \(\det X_1=-1\) en \(\FF_3\), las ecuaciones \(X_aL_a=Y_a\) determinan unívocamente \(L_a=X_a^{-1}Y_a\). En esa carta de vectores fila resultan
\[
 L_0=\begin{pmatrix}
 2&2&2&1&2&1\\2&1&2&2&1&1\\1&0&0&1&0&2\\
 0&0&1&1&1&0\\2&2&2&0&2&1\\2&1&2&1&0&0
 \end{pmatrix},\quad
 L_1=\begin{pmatrix}
 2&1&2&0&2&2\\1&2&1&1&2&0\\1&2&1&1&1&2\\
 0&1&0&0&2&1\\2&0&2&1&0&1\\0&2&2&2&1&1
 \end{pmatrix}\quad(\bmod3).
\]
El calendario \((L_0,L_0,L_1,L_1)\) produce cuatro nuevos bloques de seis componentes. Sus concatenaciones forman
\[
 w_6\longrightarrow w_{12}\longrightarrow w_{18}
 \longrightarrow w_{24}\longrightarrow w_{30}.
\]
Los subíndices de \(w\) indican longitud acumulada. La frontera \(R_{36}\) conserva la estructura terminal y abre la relación de prolongación; no se renombra como un último bloque periódico que pudiera repetirse indefinidamente.

La identidad \(L=X^{-1}Y\) demuestra unicidad una vez fijadas las transiciones; no sustituye la producción de esas transiciones por \(U_t\). La procedencia de su registro y el examen de los dieciséis calendarios binarios se documentan en el suplemento técnico. El calendario indicado es el único de esos dieciséis que reproduce conjuntamente las tres sucesiones de cinco bloques.

\subsection{Cinco regiones de clausura y compensación angular}

La región de clausura contiene una componente transversal y cuatro acompañantes. El lector simétrico actúa en el espacio de sus cinco posiciones mediante
\[
 P_5=\frac15\mathbf1\mathbf1^{\mathsf T},\qquad
 \mathbf1=(1,1,1,1,1)^{\mathsf T}.
\]
La invariancia \(P_5\lambda=\lambda\) y la conservación \(\mathbf1^{\mathsf T}\lambda=1\) fuerzan \(\lambda=\mathbf1/5\). Esta normalización distribuye la unidad entre posiciones del lector; no sustituye las multiplicidades de semillas \(432,144,144,144,144\) por frecuencias uniformes. La carta elíptica del lector realiza cada coordenada mediante \(1+\lambda_jJ\); su pendiente es \(\lambda_j\), por lo que la pendiente primitiva resulta \(1/5\). Se hace explícita así la regla que pasa de una coordenada normalizada a una pendiente, en lugar de identificar ambas nociones sin mapa. El lector compone las cuatro posiciones acompañantes por la ley
\begin{equation}
 a\oplus b=\frac{a+b}{1-ab}.
 \label{x_reciprocidad:eq:suma-pendientes}
\end{equation}
Esta ley se obtiene al multiplicar \((1+aJ)(1+bJ)\) en el régimen \(J^2=-I\) y dividir por su componente escalar. Por ello la composición de pendientes es una realización de la estructura cuadrática de TRIT, no una suma ordinaria de cuatro fracciones.

\begin{proposition}[Compensador de la clausura]
La composición de cuatro pendientes \(1/5\) es \(120/119\). La condición de retorno a pendiente uno determina un único compensador positivo de la forma \(1/q\), con \(q>1\), y da \(q=239\).
\end{proposition}
\begin{proof}
La primera duplicación da \((1/5)\oplus(1/5)=5/12\); la segunda, \((5/12)\oplus(5/12)=120/119\). La condición
\[
 \frac{120/119-1/q}{1+(120/119)/q}=1
\]
equivale a \((120-119)q=120+119\), por lo que \(q=239\).
\end{proof}

El carácter de clausura resultante tiene realización
\begin{equation}
 \Pi_{\rm cl}=4\left(4\arctan\frac15-\arctan\frac1{239}\right).
 \label{x_reciprocidad:eq:lector-pi}
\end{equation}
La identidad angular de Machin reconoce \(\Pi_{\rm cl}=\pi\). El sentido de la construcción es preciso: la pendiente y el número de composiciones están declarados por el lector de la pentafibra; la ecuación de compensación fuerza \(239\). Una vez producido este carácter, su evaluación por series alternadas no selecciona de nuevo la órbita de semillas.

Para \(0<x<1\), las sumas alternadas
\[
 A_n(x)=\sum_{j=0}^n\frac{(-1)^jx^{2j+1}}{2j+1}
\]
encierran \(\arctan x\) entre dos truncamientos consecutivos, con diferencia \(x^{2n+3}/(2n+3)\). Aplicadas a las dos pendientes de \eqref{x_reciprocidad:eq:lector-pi}, proporcionan extremos racionales de error explícito. El valor numérico se calcula así como consecuencia del carácter exacto, sin introducir un prefijo objetivo.

\subsection{Lector de propagación y normalización de la unidad}

La región de propagación se realiza mediante una familia formal \(P_t\in\mathcal A[[t]]\), donde \(\mathcal A\) es un álgebra conmutativa unitaria sobre \(\QQ\). Su composición satisface \(P_{t+s}=P_tP_s\), con unidad \(P_0=1\), y el transporte elemental orientado fija el generador escalar \(P'_0=+1\). La condición fija su signo, no sólo su magnitud. Escribiendo
\[
 P_t=\sum_{n\ge0}a_nt^n,
\]
la normalización determina \(a_0=a_1=1\). La comparación del coeficiente de \(s^nt\) en \(P_{s+t}=P_sP_t\) da \((n+1)a_{n+1}=a_na_1\); equivalentemente, \(P'_t=P_t\). Así,
\begin{equation}
 (n+1)a_{n+1}=a_n,\qquad a_n=\frac1{n!}.
 \label{x_reciprocidad:eq:propagacion}
\end{equation}
El carácter en una unidad de parámetro es \(\Pi_{\rm pr}=P_1\), reconocido posteriormente como \(e\).

La normalización del generador importa. La ley semigrupal sin ese dato admitiría \(P_t=e^{ct}\) para distintos \(c\); por ello no se omite la condición que fija la unidad del parámetro en el lector. Esta condición es anterior a la comparación decimal y pertenece a la exposición de la regla operacional, no a la elección de un valor experimental.

\begin{proposition}[Cotas racionales de propagación]
Para \(n\ge1\), sea \(E_n=\sum_{j=0}^n1/j!\). Entonces
\[
 E_n<\Pi_{\rm pr}<E_n+\frac{n+2}{(n+1)(n+1)!}.
\]
\end{proposition}
\begin{proof}
El primer término omitido es \(1/(n+1)!\). Cada cociente posterior es a lo sumo \(1/(n+2)\). La suma geométrica mayorante de la cola es
\[
 \frac1{(n+1)!}\frac1{1-1/(n+2)}
 =\frac{n+2}{(n+1)(n+1)!}.
\]
La desigualdad estricta se sigue de la disminución de los cocientes sucesivos.
\end{proof}

\subsection{Incidencia modal y generación de la autoescala}

El selector trítico produce el ciclo \((+1,-1,0)\). En los modos \(\Sigma,\Pi\), la regla es \(+1:m\mapsto\Sigma\), \(-1:m\mapsto\Pi\), \(0:m\mapsto m\). Con cierre cíclico, el modo observado es \((\Sigma,\Pi,\Pi)\). Sus aristas son \(\Sigma\to\Pi\), \(\Pi\to\Pi\) y \(\Pi\to\Sigma\). Por tanto su incidencia, en ese orden de modos, es
\begin{equation}
 F_{\rm av}=\begin{pmatrix}0&1\\1&1\end{pmatrix}.
 \label{x_reciprocidad:eq:fav}
\end{equation}
Las cuatro entradas se obtienen del calendario, no de una ecuación que haya recibido previamente el valor de \(\varphi\). La repetición del calendario da conteos \(3F_{\rm av},9F_{\rm av},36F_{\rm av}\) en longitudes nueve, veintisiete y ciento ocho. Estos conteos no son las potencias de la matriz.

\begin{theorem}[Carácter positivo de autoescala]
La incidencia \eqref{x_reciprocidad:eq:fav} posee un único rayo propio estrictamente positivo. Al normalizarlo como \((1,x)^{\mathsf T}\), su coordenada \(x\) satisface \(x^2-x-1=0\), \(x>0\), y es \(\varphi=(1+\sqrt5)/2\).
\end{theorem}
\begin{proof}
La ecuación propia da \(x=\lambda\) y \(1+x=\lambda x\). Eliminar \(\lambda\) produce el polinomio. De sus dos raíces sólo una es positiva. Como \(F_{\rm av}^2\) tiene todas sus entradas positivas, el rayo positivo es único y simple. La expresión radical reconoce la coordenada obtenida; no participa en la producción de \(F_{\rm av}\).
\end{proof}

Con \(F_0=0,F_1=1,F_{n+1}=F_n+F_{n-1}\), la incidencia induce, para \(n\geq1\),
\[
 F_{\rm av}^n=\begin{pmatrix}F_{n-1}&F_n\\F_n&F_{n+1}\end{pmatrix}.
\]
Los cocientes consecutivos \(F_{n+1}/F_n\) alternan alrededor de \(\varphi\). La identidad de determinante
\[
 F_{n+1}^2-F_nF_{n+2}=(-1)^n
\]
da una anchura exacta \(1/(F_nF_{n+1})\) entre dos cocientes consecutivos. Su crecimiento hace arbitrariamente pequeña esta horquilla racional.

La medida sobre todas las historias compatibles no coincide con la frecuencia de modos del calendario de tres eventos. La frecuencia cronológica es \((1/3,2/3)\). La medida de máxima entropía del envolvente simbólico definido por \(F_{\rm av}\) se construye con sus vectores propios y tiene razón de pesos \(\varphi^{-2}\). Esta última interviene en la lectura areal--volumétrica; no se obtiene sustituyendo sin más una frecuencia de conteo por otra.


% BEGIN OWNER /Users/ruben/Documents/New project/output/EXTREMA_MEDIA_RAZON_RECIPROCIDAD_ES_EN_20260918/ARTICULO/ES/source/sections/retorno_areal_volumetrico.tex
\subsubsection{Espacio de trayectorias y medida invariante entre hojas}
\label{x_reciprocidad:sec:medida-retorno-hojas}

La topología toroidal de APP fija el soporte de las traslaciones cardinales
y sus clases de enrollamiento. La correspondencia entre hojas procede de
otra operación sobre ese mismo soporte: la selección temporal de la
evaluación aditiva o multiplicativa. Su tipado geométrico es
\[
 \eta_{\rm av}(\Sigma)=\mathscr H_{\rm ar},\qquad
 \eta_{\rm av}(\Pi)=\mathscr H_{\rm vol}.
\]
La primera hoja corresponde a la lectura aditiva de frontera; la segunda,
a la lectura multiplicativa de interior. Este mapa transporta las etiquetas
de modo y conserva la incidencia \eqref{x_reciprocidad:eq:fav}. El cuadrado de la razón
de autoescala se determina mediante las trayectorias de esa incidencia.

\Needspace{8\baselineskip}
\begin{definition}[Envolvente simbólico de memoria uno]
Sea
\[
 X_{\rm av}=\{(s_k)_{k\in\mathbb Z}\in
 \{\mathscr H_{\rm ar},\mathscr H_{\rm vol}\}^{\mathbb Z}:
 (F_{\rm av})_{s_k,s_{k+1}}=1\text{ para todo }k\}.
\]
El desplazamiento de índices actúa sobre este espacio. Es el menor sistema
simbólico de memoria uno que contiene las sucesiones modales y conserva
sus pares consecutivos: admite las dos transiciones cruzadas y la
autorretención volumétrica. La fase, la orientación cardinal y los registros
de la trayectoria permanecen en el estado TPK del que se obtiene este
envolvente.
\end{definition}

\begin{proposition}[Ponderación de las trayectorias entre hojas]
En el orden areal--volumétrico, la medida invariante de entropía máxima
de \(X_{\rm av}\) tiene matriz de transición y distribución estacionaria
\begin{equation}
 P_{\rm av}=\begin{pmatrix}0&1\\\varphi^{-2}&\varphi^{-1}\end{pmatrix},
 \qquad
 (\mu_{\rm ar},\mu_{\rm vol})
 =\frac{(1,\varphi^2)}{1+\varphi^2}.
 \label{x_reciprocidad:eq:medida-hojas-exacta}
\end{equation}
En particular, \(\mu_{\rm ar}/\mu_{\rm vol}=\varphi^{-2}\).
\end{proposition}
\begin{proof}
La incidencia ya construida tiene vector propio positivo
\(r=(1,\varphi)^{\mathsf T}\) y autovalor \(\varphi\). La
normalización
\[
 (P_{\rm av})_{ij}
 =\frac{(F_{\rm av})_{ij}r_j}{\varphi r_i}
\]
da la matriz indicada. Sus filas suman uno por
\(\varphi^{-2}+\varphi^{-1}=1\). La simetría de \(F_{\rm av}\)
implica que los pesos proporcionales a \(r_i^2\) satisfacen balance
detallado; por tanto, \(\mu P_{\rm av}=\mu\). La irreducibilidad y la
autorretención volumétrica determinan una única distribución estacionaria.

Para verificar la maximalidad, en toda transición permitida se cumple
\(\log(P_{\rm av})_{ij}=\log r_j-\log r_i-\log\varphi\).
Bajo cualquier medida invariante, los términos de extremo se cancelan
en promedio. Si \(\nu_i\) son sus pesos de vértice y
\(q_i(j)=\nu(s_1=j\mid s_0=i)\), para \(\nu_i>0\), entonces
\[
 h_\nu\leq H_\nu(s_1\mid s_0)
 =\log\varphi-\sum_{i:\nu_i>0}\nu_iD(q_i\Vert P_i)
 \leq\log\varphi,
\]
donde \(D(q\Vert p)=\sum_jq_j\log(q_j/p_j)\) es la entropía
relativa en los pares permitidos y se adopta \(0\log0=0\).
La medida markoviana de \eqref{x_reciprocidad:eq:medida-hojas-exacta} alcanza la cota.
La igualdad en la primera desigualdad exige la propiedad de Markov;
la igualdad en la segunda exige \(q_i=P_i\). La estacionariedad y
la irreducibilidad fijan entonces \(\mu\), lo que establece la
unicidad. Esta es la medida de Parry del envolvente definido.
\end{proof}

La frecuencia \((1/3,2/3)\) cuenta los instantes del calendario primitivo;
\(\mu\) pondera las trayectorias del envolvente. Su razón cuadrática
expresa el producto de las amplitudes de incidencia de entrada y salida.
La memoria y las restricciones de fase del TPK conservan su dominio
propio; la entropía del envolvente corresponde al sistema simbólico
especificado en esta definición.

\subsubsection{Operador de clausura y límite de los retornos marcados}
\label{x_reciprocidad:sec:operador-retorno-marcado}

La coordenada de clausura \(\pi\), generada por el lector regional,
interviene después de la construcción de la incidencia modal. Sean
\(e_{\rm ar}=(1,0)^{\mathsf T}\),
\(e_{\rm vol}=(0,1)^{\mathsf T}\) y
\(P_i=e_ie_i^{\mathsf T}\) los proyectores de hoja. Las condiciones
del lector son conservación de ambas hojas, normalización volumétrica
unitaria y marca areal igual a la coordenada de clausura. Estas condiciones
determinan el operador positivo
\begin{equation}
 D_\pi=\pi P_{\rm ar}+P_{\rm vol}
       =\begin{pmatrix}\pi&0\\0&1\end{pmatrix}.
 \label{x_reciprocidad:eq:operador-marca-pi}
\end{equation}
La diagonalidad expresa la conservación de hojas y las dos entradas
expresan sus normalizaciones; de este modo queda explícita la unicidad
del operador para el lector fijado.

\Needspace{8\baselineskip}
\begin{proposition}[Razón de retornos de frontera]
Para \(n\geq1\), la evaluación marcada
\begin{equation}
 \Theta_n=
 \frac{\langle e_{\rm ar},D_\pi F_{\rm av}^{\,n}e_{\rm ar}\rangle}
 {\langle e_{\rm vol},F_{\rm av}^{\,n}e_{\rm vol}\rangle}
 =\pi\frac{F_{n-1}}{F_{n+1}}
 \label{x_reciprocidad:eq:retorno-marcado-finito}
\end{equation}
satisface
\begin{equation}
 \lim_{n\to\infty}\Theta_n
 =\pi\frac{\mu_{\rm ar}}{\mu_{\rm vol}}
 =\frac{\pi}{\varphi^2}.
 \label{x_reciprocidad:eq:retorno-marcado-limite}
\end{equation}
\end{proposition}
\begin{proof}
Las entradas diagonales de \(F_{\rm av}^{\,n}\) cuentan los caminos
de longitud \(n\) que regresan a su hoja inicial. La fórmula de potencias
ya demostrada proporciona \(F_{n-1}\) y \(F_{n+1}\). El operador
\(D_\pi\) multiplica por \(\pi\) el primer recuento. Finalmente,
\(F_{n-1}/F_{n+1}\) es el producto de dos cocientes consecutivos
que convergen a \(\varphi^{-1}\); su límite es
\(\varphi^{-2}\). La identidad con la razón de pesos se sigue de
\eqref{x_reciprocidad:eq:medida-hojas-exacta}.
\end{proof}

El operador de clausura aparece una sola vez, en la evaluación terminal
del retorno. La transferencia \(F_{\rm av}\) permanece determinada
por el calendario. Así se compone la información de dos lectores del
mismo estado: la incidencia modal determina la ponderación cuadrática
y la región de clausura aporta su coordenada a la frontera.

La composición conserva también el refinamiento. Para cilindros positivos
\(I=[a,b]\) y \(J=[c,d]\), la imagen del lector
\(\mathcal R(x,y)=x/y^2\) es el intervalo
\[
 \mathcal Q(I,J)=\left[\frac{a}{d^2},\frac{b}{c^2}\right].
\]
La monotonía en cada argumento prueba que los cilindros encajados de
clausura y autoescala producen intervalos encajados de la razón; la
positividad de sus límites y la contracción de sus diámetros determinan
\(\pi/\varphi^2\). La expresión escalar retiene así una cadena
de operaciones compatible con la profundidad, que se aplica a
continuación en sus dos orientaciones.

% END OWNER


Después de generados los caracteres, las dos orientaciones del lector
\[
 R(x,y)=\frac{x}{y^2}
\]
proporcionan
\begin{equation}
 \rho_A=\frac\pi{\varphi^2},\qquad
 \rho_E=\frac\varphi{\pi^2},\qquad
 \varphi^3\rho_A^2\rho_E=1.
 \label{x_reciprocidad:eq:areal-electron}
\end{equation}
La segunda expresión interviene en el exponente de la construcción electrónica de la sección~\ref{x_reciprocidad:el-centro:registros}. La primera se obtiene también como límite de la lectura marcada \(\pi F_{n-1}/F_{n+1}\). El intercambio de argumentos relaciona ambas, pero no convierte una reparametrización de la ecuación electrónica en una segunda derivación independiente de todos sus coeficientes.


% BEGIN OWNER /Users/ruben/Documents/New project/output/EXTREMA_MEDIA_RAZON_RECIPROCIDAD_ES_EN_20260918/ARTICULO/ES/source/sections/euler_trit_articulacion.tex
% Articulación posterior a la generación de pi, phi y e.
% Contenido acumulativo; no reemplaza el desarrollo previo del TRIT.
\subsection{Realizaciones exponenciales de las coordenadas generadas}
\label{x_reciprocidad:euler-trit-articulacion}

Las coordenadas regionales ya determinadas permiten reconocer las
realizaciones concretas de la familia cuadrática del TRIT. El orden es
sustantivo: la construcción de las regiones precede a la evaluación de una
exponencial en los parámetros \(\pi\) o \(2\log\varphi\). Estas expresiones
identifican la función de las coordenadas en operadores construidos desde
APP y TPK; no intervienen en la selección de los prefijos que las producen.

\subsubsection{Ciclo ternario, transporte nilpotente e incidencia de hojas}

La multiplicación \(a(r)=4r\) en \(\mathbb Z/9\mathbb Z\) tiene orden
\(3\). Sobre las unidades determina las órbitas \(\{1,4,7\}\) y
\(\{2,8,5\}\). En el módulo de aumento de cualquiera de ellas, formado por
las combinaciones enteras cuyos coeficientes suman cero, la base
\((e_r-e_{4r},e_{4r}-e_{7r})\) da
\[
 C=\begin{pmatrix}0&-1\\1&-1\end{pmatrix},
 \qquad C^2+C+I=0.
\]
La forma restringida tiene matriz de Gram
\(\bigl(\begin{smallmatrix}2&-1\\-1&2\end{smallmatrix}\bigr)\).
En las evaluaciones residuales módulo \(9\), la acción distingue las hojas:
\(S(ax,ay)=aS(x,y)\), mientras \(P(ax,ay)=a^{-1}P(x,y)\).
La realización cíclica conserva, por tanto, una orientación contragrediente
entre suma y producto. Los cocientes del levantamiento entero se conservan
en el registro aritmético correspondiente.

El transporte parabólico elemental se representa mediante
\[
 N=\begin{pmatrix}0&1\\0&0\end{pmatrix}.
\]
La incidencia areal--volumétrica ya construida en \eqref{x_reciprocidad:eq:fav} proporciona
\[
 F_{\rm av}=\begin{pmatrix}0&1\\1&1\end{pmatrix},
 \qquad A=F_{\rm av}^2=\begin{pmatrix}1&1\\1&2\end{pmatrix}.
\]
Los tres operadores actúan en espacios reales declarados: el módulo de
aumento tensorizado con \(\mathbb R\), el plano del transporte nilpotente y
el plano de incidencia modal. Compartir dimensión matricial no identifica
estos espacios ni sus coordenadas.

\begin{proposition}[Generadores concretos de los tres tipos]
\label{x_reciprocidad:euler-trit-generadores}
Los operadores
\[
 J_{\mathrm e}=\frac{2C+I}{\sqrt3},\qquad
 J_{\mathrm p}=N,\qquad
 J_{\mathrm h}=\frac{2A-3I}{\sqrt5}
\]
satisfacen, respectivamente,
\[
 J_{\mathrm e}^{\,2}=-I,\qquad
 J_{\mathrm p}^{\,2}=0,\qquad
 J_{\mathrm h}^{\,2}=I.
\]
\end{proposition}
\begin{proof}
La relación de \(C\) implica \((2C+I)^2=-3I\).
La multiplicación directa da \(N^2=0\).
Finalmente, \(A\) tiene traza \(3\) y determinante \(1\), por lo que
\(A^2-3A+I=0\) y \((2A-3I)^2=5I\).
\end{proof}

\begin{proposition}[Tres evaluaciones de la identidad exponencial]
\label{x_reciprocidad:euler-trit-evaluaciones}
En las realizaciones anteriores,
\begin{equation}
 \boxed{
 \exp(\pi J_{\mathrm e})+I=0,\qquad
 \exp(J_{\mathrm p})=I+J_{\mathrm p},\qquad
 \exp(2\log\varphi\,J_{\mathrm h})=F_{\rm av}^{\,2}.}
 \label{x_reciprocidad:euler-trit-tres-evaluaciones}
\end{equation}
\end{proposition}
\begin{proof}
La especialización elíptica de la exponencial trítica da
\(\cos\pi\,I+\sin\pi\,J_{\mathrm e}=-I\).
La nilpotencia elimina exactamente los términos de grado al menos \(2\)
en la segunda expresión. Para la tercera, \(\varphi^2=\varphi+1\) implica
\[
 \cosh(2\log\varphi)=\frac32,\qquad
 \sinh(2\log\varphi)=\frac{\sqrt5}{2}.
\]
Por tanto, la exponencial vale
\(\frac32I+\frac{\sqrt5}{2}J_{\mathrm h}=A\).
\end{proof}

La primera igualdad representa la media vuelta elíptica; la vuelta completa
satisface \(\exp(2\pi J_{\mathrm e})=I\). La segunda mantiene un término
lineal no nulo: el régimen parabólico expresa transporte, no ausencia de
evolución. La tercera realiza un avance hiperbólico unimodular. Sus
autovalores son \(\varphi^2\) y \(\varphi^{-2}\); una dirección se expande y
la otra se contrae conservando el determinante. Las tres evaluaciones
emplean parámetros distintos de una familia exponencial común.

En esa familia, \(e\) interviene mediante la evaluación escalar
\(\exp(I)=eI\) y la ley \(\exp(sI)\exp(tI)=\exp((s+t)I)\).
Su papel no se restringe al tipo parabólico. La coordenada \(e\), obtenida
previamente de su región, es reconocida aquí como la evaluación unitaria
de la exponencial; esta identidad no sustituye su genealogía coinductiva.

\subsubsection{Resolución polinómica de los regímenes}

Las tres realizaciones admiten una presentación compacta sobre la suma
directa de sus espacios. Defínanse
\[
 \mathbb J=J_{\mathrm e}\oplus J_{\mathrm p}\oplus J_{\mathrm h},
 \qquad Q=\mathbb J^2.
\]
Las relaciones anteriores implican
\[
 \mathbb J^6=\mathbb J^2,\qquad Q^3=Q.
\]
El uso de \(Q\) mantiene separado este cuadrado operatorial del registro
dodecafásico \(K\).

\begin{proposition}[Proyectores polinómicos de régimen]
\label{x_reciprocidad:euler-trit-proyectores}
Las aplicaciones
\[
 p_{\mathrm e}=\frac{Q^2-Q}{2},\qquad
 p_{\mathrm p}=I-Q^2,\qquad
 p_{\mathrm h}=\frac{Q^2+Q}{2}
\]
son idempotentes mutuamente ortogonales, suman \(I\) y recuperan los tres
sumandos de la representación.
\end{proposition}
\begin{proof}
Los tres polinomios toman los valores \((1,0,0)\), \((0,1,0)\) y
\((0,0,1)\) en los puntos \(-1,0,+1\), respectivamente. Las diferencias
entre sus cuadrados y ellos mismos, así como sus productos cruzados, son
múltiplos de \(X^3-X\). Sustituir \(Q\), que anula ese polinomio, demuestra
las identidades. En los tres bloques, \(Q\) actúa como \(-I,0,I\).
\end{proof}

Se obtiene así
\begin{align}
 \exp(t\mathbb J)
 ={}&p_{\mathrm e}(\cos t\,I+\sin t\,\mathbb J)
   +p_{\mathrm p}(I+t\mathbb J)\notag\\
   &+p_{\mathrm h}(\cosh t\,I+\sinh t\,\mathbb J).
 \label{x_reciprocidad:euler-trit-exponencial-total}
\end{align}
Cada sumando conserva su relación cuadrática y la conjugación invierte
su evolución. La representación total reúne los tres regímenes, pero no
define por sí sola un operador de transición entre ellos. Tal transición
requiere los mapas de transporte del TPK con sus dominios, orientación y
memoria. La articulación exponencial permite reconocer sus tipos sin
reemplazar esa dinámica por una identificación de las fibras.

% PROCEDENCIA FOCAL (no impresa)
% Resultado recuperado: tres generadores y tres evaluaciones exponenciales.
% Propietario completo:
% BASE/manuscrito/sections/hmt/09b_algebras_cuadraticas.tex:289-516.
% Antecedente causal de C y acción contragrediente:
% BASE/manuscrito/sections/hmt/03d_esqueleto_ciclotomico_app.tex:13-47,145-246.
% Los tres tipos y la orientación están en:
% BASE/manuscrito/sections/hmt/02_trit_geometrias.tex:83-196.
% Resolución polinómica recuperada:
% output/INVESTIGACION_HOLONOMIA_NONADICA_CRISTAL_APERIODICO_20260903/
% ALGEBRA_HOLONOMICA_UNIVERSAL_Y_EULER_TRITICO.md, secciones 8-11 y 22.
% Los cuatro propietarios se leyeron completos. BASE designa:
% /Users/ruben/Documents/HMT2/
% EL_CIERRE_HOLOGRAFICO_DEL_INFINITO_SUCESOR_102_CAPITULOS_2026-08-28_EN_TRABAJO/
% No se incorpora la identidad P9(alpha)=delta del documento especializado:
% esa identificación no pertenece a este bloque.
% Tampoco se identifica una dirección nilpotente con la vacancia electrónica,
% ni se presume una representación global del transporte por la suma directa.
% COMPROBACIONES: Python estándar, enteros y Fraction.
% C^2+C+I=0; (2C+I)^2=-3I; N^2=0; (2Fav^2-3I)^2=5I.
% 81 pares de residuos verifican las dos acciones contragredientes.
% Los 3 proyectores se verificaron en Q=-1,0,+1; prueba general en texto.
% 6 labels únicos y entornos equilibrados. Sin compilación.

% END OWNER


\subsection{Cilindros compatibles y profundidad arbitraria}

La realización de un bloque ternario \(b=(b_1,\ldots,b_6)\) es el entero
\[
 \nu(b)=\sum_{j=1}^6b_j3^{6-j}\in\{0,\ldots,728\}.
\]
Una sucesión de bloques define prefijos enteros por \(N_{n+1}=729N_n+\nu(b_{n+1})\), comenzando en \(N_0=m\), donde \(m\) es la parte entera determinada por el lector. Las horquillas anteriores sitúan los caracteres de clausura, propagación y autoescala en \((3,4)\), \((2,3)\) y \((1,2)\), respectivamente; por tanto, sus valores iniciales son \(m=3,2,1\). Los bloques siguientes refinan la parte fraccionaria. El cilindro correspondiente es
\begin{equation}
 I_n=\left[\frac{N_n}{729^n},\frac{N_n+1}{729^n}\right],
 \qquad |I_n|=729^{-n}.
 \label{x_reciprocidad:eq:cilindro}
\end{equation}
La convención semiabierta se utiliza para seleccionar un representante cuando un valor cae en una frontera. Para el paso al límite se utilizan las clausuras de estos intervalos y se identifican las dos expansiones fronterizas equivalentes. Esta precisión impide suponer que toda intersección de intervalos semiabiertos encajados sea automáticamente no vacía.

\begin{theorem}[Determinación de la realización escalar]
Una historia de bloques compatible determina un único límite de los extremos izquierdos de \eqref{x_reciprocidad:eq:cilindro}. Si un carácter exacto dispone de horquillas racionales de anchura arbitrariamente pequeña y se separa de la frontera de cada cilindro decimal examinado, cada prefijo decimal de longitud finita se determina después de un cálculo finito. Los prefijos así obtenidos son compatibles entre sí.
\end{theorem}
\begin{proof}
La extensión de un bloque conserva el prefijo anterior por división euclídea. Los cilindros cerrados están encajados y su diámetro tiende a cero; sus extremos izquierdos son de Cauchy y tienen un único límite. Para una longitud decimal fijada, un valor separado de las fronteras de la partición posee distancia positiva a ellas. Una horquilla suficientemente estrecha queda contenida en una única celda y determina su prefijo. La unicidad de la celda y el encajamiento de las particiones prueban la compatibilidad por truncamiento. En los puntos con expansión terminal debe añadirse la decisión exacta de frontera y su convención de representación; no se deduce esa decisión de una anchura numérica solamente.
\end{proof}

\begin{corollary}[Obtención finita y compatibilidad de los prefijos regionales]
\label{x_reciprocidad:gen:prefijos-regionales}
Para cada una de las realizaciones \(\pi,\varphi,e\) y para toda longitud
decimal finita, las horquillas construidas determinan el prefijo correspondiente
mediante un cálculo finito. Los prefijos obtenidos a distintas profundidades
son compatibles por truncamiento.
\end{corollary}
\begin{proof}
Los caracteres de clausura, propagación y autoescala poseen las horquillas
anteriores. Tras su reconocimiento, la irracionalidad de \(\pi,e,\varphi\)
garantiza su separación de las fronteras racionales de cada partición decimal
finita. Se aplica el teorema precedente a cada carácter.
\end{proof}
La profundidad arbitraria expresa el cuantificador del corolario: para todo
horizonte finito existe un cálculo que determina su prefijo. El objeto
matemático es la familia compatible completa; cada ejecución obtiene una
proyección finita de ella.

\subsection{Involución especular y sincronización de las regiones}

La conexión nonádica transporta conjuntamente los tres caracteres. En la carta de la novena fase, el eje de autoescala y la suma ternaria de clausura y propagación permanecen fijados bajo el intercambio especular. La inversión helicoidal cambia el sentido del transporte; el espejo de los dos canales cambia su orientación relativa. Son involuciones diferentes, aunque ambas actúen sobre el mismo sistema de trayectorias.

Las coincidencias de censo ilustran esa coordinación:
\begin{center}
\begin{tabular}{@{}c rrr l@{}}\toprule
Fase&\(N_\pi\)&\(N_e\)&\(N_\varphi\)&Igualdad de conteos\\\midrule
4&207&207&122&\(N_\pi=N_e\)\\
5&39&151&151&\(N_e=N_\varphi\)\\
6&110&110&69&\(N_\pi=N_e\)\\
7&81&29&81&\(N_\varphi=N_\pi\)\\
8&59&59&59&sincronización triple\\
9&43&43&43&sincronización triple\\\bottomrule
\end{tabular}
\end{center}
Estos enteros cuentan coordenadas compatibles de las regiones. No afirman una igualdad entre \(\pi\), \(e\) y \(\varphi\) como números reales. Si
\[
 \partial(a,b,c)=(a-b,b-c,c-a),
\]
su imagen pertenece al retículo \(A_2=\{(u,v,w)\in\ZZ^3:u+v+w=0\}\). En las fases ocho y nueve la componente relativa se anula, pero la componente diagonal cambia de \((59,59,59)\) a \((43,43,43)\). La sincronización relativa coexiste con un desplazamiento de \(-16\) por canal. El cierre que conducirá a \(\alpha\) debe conservar ambas informaciones: coincidencia relativa y evolución de la memoria.

\subsection{Suspensión helicoidal esturmiana}

La relación entre las bases nueve y diez de la completación define el parámetro adimensional
\begin{equation}
 \delta=\frac{\log(10/9)}{\log10}.
 \label{x_reciprocidad:eq:delta-sturm}
\end{equation}
No se obtiene de una constante física objetivo. Es la escala logarítmica relativa de dos bases ya presentes en la construcción. Su empleo como lectura simbólica del refinamiento tiene una consecuencia exacta.

\begin{lemma}[Irracionalidad de la pendiente]
El número \(\delta\) es irracional.
\end{lemma}
\begin{proof}
Si \(\delta=p/q\), \(q>0\), entonces \((10/9)^q=10^p\) y \(9^q=10^{q-p}\). La factorización del primer miembro contiene el primo tres y la del segundo sólo los primos dos y cinco. La igualdad es imposible.
\end{proof}

La palabra mecánica inferior, con intercepto \(\theta\), es
\[
 z_n(\theta)=\lfloor(n+1)\delta+\theta\rfloor-\lfloor n\delta+\theta\rfloor\in\{0,1\}.
\]
Se obtiene al observar la rotación \(\theta\mapsto\theta+\delta\pmod1\) mediante el intervalo \([1-\delta,1)\). La sucesión de eventos tiene densidad \(\delta\), y la suma sobre cualquier ventana de longitud \(L\) vale \(\lfloor L\delta\rfloor\) o \(\lceil L\delta\rceil\). Como \(21<1/\delta<22\), las separaciones entre eventos son veintiuno o veintidós. La palabra que indica cuáles son los intervalos cortos es otra palabra mecánica de pendiente \(\sigma=22-1/\delta\); sus eventos están separados por seis o siete intervalos primarios. La segunda escala es inducida por la primera y no un segundo parámetro libre.


% BEGIN OWNER /Users/ruben/Documents/New project/output/EXTREMA_MEDIA_RAZON_RECIPROCIDAD_ES_EN_20260918/ARTICULO/ES/source/sections/vacancias_capacidad.tex
% Desarrollo reunido de DESARROLLO_MATEMATICO, apartados 38--40 y 50.
\subsubsection{Índice de capacidad, vacancias e inducción del régimen trítico}
\label{x_reciprocidad:vac-capacidad-trit}

La sucesión de vacancias está vinculada al refinamiento mediante una
identidad de incrementos. Para mostrarla se fija el origen de fase
\(\theta=0\). Sean
\[
 \rho=\log_{10}9=1-\delta,\qquad
 \mathcal C_t=\lfloor(t+1)\rho\rfloor\quad(t\geq0).
\]
El índice \(\mathcal C_t\) expresa la capacidad decimal del refinamiento
en la normalización que compara un bloque ternario de \(6\) posiciones,
con \(729\) posibilidades, y un bloque decimal de \(3\) posiciones, con
\(1000\) posibilidades. En efecto,
\(\log_{1000}729=\log_{10}9=\rho\). Se reserva la letra \(K\) para
el registro dodecafásico que se construirá posteriormente.

\begin{proposition}[Complementariedad de capacidad y vacancia]
\label{x_reciprocidad:prop:capacidad-vacancia}
Para \(t\geq1\),
\begin{equation}
 \mathcal C_t-\mathcal C_{t-1}=1-z_t(0).
 \label{x_reciprocidad:eq:capacidad-vacancia}
\end{equation}
En consecuencia, si \(0\leq a<b\),
\[
 \mathcal C_b-\mathcal C_a+
 \sum_{t=a+1}^{b}z_t(0)=b-a.
\]
\end{proposition}
\begin{proof}
La irracionalidad de \(\delta\) da
\[
 \mathcal C_t
 =t+1-\lceil(t+1)\delta\rceil
 =t-\lfloor(t+1)\delta\rfloor.
\]
La resta de dos términos consecutivos produce
\eqref{x_reciprocidad:eq:capacidad-vacancia}. La suma sobre el intervalo es telescópica.
\end{proof}

Cada transición aumenta en una unidad el contador de bloques. Si
\(z_t=0\), aumenta también la capacidad; si \(z_t=1\), la capacidad se
mantiene mientras se incorpora una transición a la historia. La vacancia
es aquí un evento de retención de capacidad, con una posición y un
antecedente determinados. El balance anterior conserva exactamente la
longitud entre capacidad adquirida y vacancias registradas. La fase
observada \(g_t^{\rm cap}=[\mathcal C_t]_9\) satisface
\[
 g_t^{\rm cap}-g_{t-1}^{\rm cap}\equiv1-z_t\pmod9.
\]
Esta fase de capacidad y la fase cronológica \([t]_9\) son dos
observaciones diferentes. Una retención local de capacidad y una vuelta
completa de la holonomía tienen mecanismos distintos, aunque ambas
permitan conservar la observación de fase con modificación del estado.

La relación con TRIT puede seguirse en las propias posiciones de vacancia.
Para \(j\geq1\), defínanse
\[
 p_j=\left\lfloor\frac j\delta\right\rfloor,\qquad
 v_j=\mathcal C_{p_j},\qquad h_j=p_{j+1}-p_j.
\]
La desigualdad \(p_j\delta<j<(p_j+1)\delta\) identifica \(p_j\)
como la posición del evento \(j\). Como \(0<\delta<1\), da también
\(\lfloor(p_j+1)\delta\rfloor=j\), y por tanto
\[
 v_j=p_j-j,\qquad
 v_{j+1}-v_j=h_j-1.
\]
Si \(s_j=22-h_j\), entonces \(s_j\in\{0,1\}\) y
\begin{equation}
 [v_{j+1}]_3=[v_j-s_j]_3.
 \label{x_reciprocidad:eq:vacancia-regimen}
\end{equation}
Un intervalo de longitud \(22\) conserva la clase; uno de longitud
\(21\) la desplaza una unidad en sentido negativo. El selector canónico
se aplica a la fase positiva \(r_j=1+[v_j]_9\):
\[
 \tau_j=M_{\rm ph}(r_j).
\]
Queda determinado así un régimen trítico por el índice de capacidad
retenido en cada evento, además de la codificación binaria de la vacancia.

Las posiciones de los intervalos cortos forman la codificación mecánica
inducida de pendiente \(22-1/\delta\). Sus separaciones \(6\) y \(7\)
determinan las longitudes de las mesetas de la clase \([v_j]_3\), una
vez fijados el origen y la convención de extremos. Las primeras clases
son \(2\) durante \(6\) eventos, \(1\) durante \(7\) y \(0\) durante
\(7\); el primer tramo de \(6\) es la restricción inicial de una meseta.
Sus firmas recorren \(0,-1,+1\). Las ecuaciones
\eqref{x_reciprocidad:eq:capacidad-vacancia} y \eqref{x_reciprocidad:eq:vacancia-regimen} explicitan el
enlace entre refinamiento, vacancias y selección del régimen. La lectura
cuadrática del TRIT realiza después esos tipos mediante sus respectivos
generadores; el acontecimiento temporal y su representación geométrica
conservan así una dependencia identificable.

La inversión del índice monótono proporciona otra observación útil:
\[
 \mathcal T(k)=\min\{t:\mathcal C_t=k\}
   =\left\lceil\frac k\rho\right\rceil-1\qquad(k\geq1).
\]
Si \(\sigma=\rho^{-1}-1\), el número de bloques necesarios para ganar
\(m\) unidades de capacidad es
\[
 \mathcal T(k+m)-\mathcal T(k)
 =m+\lceil(k+m)\sigma\rceil-\lceil k\sigma\rceil
 \in\{m+\lfloor m\sigma\rfloor,m+\lceil m\sigma\rceil\}.
\]
En particular, \(9\) unidades de capacidad requieren \(9\) o \(10\)
bloques. El período de una fase finita permanece exacto, mientras el
tiempo de retorno medido por el otro índice admite dos valores. Esta
distinción es parte de la dinámica del refinamiento y no una corrección
de sus valores numéricos a posteriori.

% END OWNER


\begin{theorem}[Complejidad de las lecturas de fase]
Para \(m\ge1\), la palabra mecánica tiene complejidad factorial \(p(m)=m+1\). Al conservar la fase \(n\bmod9\), la complejidad es \(p_9(m)=9(m+1)\); al conservar sólo su clase ternaria, es \(p_3(m)=3(m+1)\). Las tres entropías topológicas son nulas.
\end{theorem}
\begin{proof}
Los predecesores de los dos extremos del intervalo de codificación forman, para factores de longitud \(m\), \(m+1\) puntos distintos de la circunferencia: la coincidencia entre extremos desplazados elimina las duplicaciones y la irracionalidad impide coincidencias adicionales. Esos puntos determinan \(m+1\) intervalos de itinerario constante y distintos. Toda órbita irracional visita cada intervalo, de donde \(p(m)=m+1\).

Fijada una fase módulo nueve, las posiciones de esa fase recorren una rotación de paso \(9\delta\), también irracional. Por densidad aparecen todos los factores en cada fase, y la primera coordenada distingue sus nueve copias. El mismo argumento con paso \(3\delta\) da tres copias tras la reducción ternaria. Finalmente, \(m^{-1}\log(c(m+1))\to0\) para \(c=1,3,9\).
\end{proof}

El producto de fases \(C_9\times\mathbb T\), trasladado por \((1,\delta)\), tiene caracteres de frecuencia
\[
 \frac a9+k\delta\pmod1,\qquad a\in\ZZ/9\ZZ,\ k\in\ZZ.
\]
El módulo espectral no es \(\log10\): es el módulo aditivo generado por \(1/9\) y \(\delta\), módulo los enteros. El factor \(\log10\) pertenece a la normalización de \eqref{x_reciprocidad:eq:delta-sturm}.

Una presentación operatoria hace explícita la memoria. En \(\ell^2(\ZZ)\), sea \(T|n\rangle=|n+1\rangle\), \(V_9|n\rangle=\zeta_9^n|n\rangle\), \(W_\delta|n\rangle=e^{2\pi i n\delta}|n\rangle\) y \(D_M|n\rangle=\lfloor n/9\rfloor|n\rangle\). Sobre vectores de soporte finito,
\begin{equation}
 V_9T=\zeta_9TV_9,\qquad
 W_\delta T=e^{2\pi i\delta}TW_\delta,\qquad
 [D_M,T^9]=T^9.
 \label{x_reciprocidad:eq:weyl-relojes}
\end{equation}
Las dos primeras relaciones son covariancias de Weyl; la tercera expresa que una vuelta nonádica aumenta el grado de memoria. La rotación de base es abeliana, mientras el álgebra de observables y traslaciones no lo es. La inversión \(|n\rangle\mapsto|-n\rangle\) conjuga el avance con el retroceso. Esta realización proporciona una suspensión helicoidal aperiódica con retorno finito de fase, no una palabra periódica de longitud nueve.


% BEGIN OWNER /Users/ruben/Documents/New project/output/EXTREMA_MEDIA_RAZON_RECIPROCIDAD_ES_EN_20260918/ARTICULO/ES/source/figures/holonomia.tex
\begin{figure}[htbp]
\centering
\begin{tikzpicture}[x=1cm,y=1cm,>=Latex,font=\small]
\begin{scope}[xshift=1.8cm,yshift=.3cm]
\draw[->,gray] (0,-.5)--(0,3.65) node[above] {\(n/9\)};
\foreach \k in {0,1,2,3}{
 \draw[gray!60,dashed] (-1.1,\k)--(1.1,\k);
 \node[left,font=\scriptsize] at (-1.1,\k){\(\k\)};
}
\draw[black,line width=.65pt,domain=0:27,samples=260,smooth,variable=\n]
 plot ({cos(40*\n)},{\n/9+.14*sin(40*\n)});
\foreach \n in {0,...,27}{
 \fill ({cos(40*\n)},{\n/9+.14*sin(40*\n)}) circle (.026);
}
\foreach \k in {0,1,2,3}{\fill (1,\k) circle (.045);}
\node[align=center,text width=3.3cm,font=\scriptsize] at (0,-1.02)
 {Misma fase en \(n=0,9,18,27\);\\grado de memoria distinto.};
\end{scope}
\begin{scope}[xshift=4.2cm,yshift=.5cm]
\node[anchor=west] at (0,3.1){\(z_n(0)=\lfloor(n+1)\delta\rfloor-\lfloor n\delta\rfloor\)};
\draw[->] (0,1.7)--(9,1.7) node[right]{\(n\)};
\foreach \n in {21,43,65,87,109,131,152,174,196,218}{
 \draw[line width=.6pt] ({\n*.039},1.7)--({\n*.039},2.25);
 \fill ({\n*.039},2.25) circle (.033);
 \node[below,font=\tiny] at ({\n*.039},1.7){\n};
}
\foreach \a/\b/\g in {21/43/22,43/65/22,65/87/22,87/109/22,109/131/22,131/152/21,152/174/22,174/196/22,196/218/22}{
 \node[font=\scriptsize] at ({(\a+\b)*.0195},2.55){\g};
}
\node[anchor=west,align=left,text width=8.3cm] at (0,.55)
 {La codificación irracional determina los intervalos entre eventos. El retorno de la fase nonádica no convierte esta secuencia en una palabra periódica.};
\end{scope}
\end{tikzpicture}
\caption{Representaciones complementarias del transporte. A la izquierda, proyección de la hélice \((\cos(2\pi n/9),\sin(2\pi n/9),n/9)\): el avance de nueve pasos recupera la fase e incrementa la altura. A la derecha, los diez primeros eventos de la palabra inferior para \(\delta=\log_{10}(10/9)\), con separaciones exactas \(21\) o \(22\). La representación helicoidal no introduce una métrica física adicional.}
\label{x_reciprocidad:fig:holonomia-esturmiana}
\end{figure}

% END OWNER


Para este recuento fijamos el intercepto \(\theta=0\) y los extremos acumulados \((N_0,\ldots,N_4)=(0,729,972,999,1000)\), correspondientes al orden cúbico \((729,243,27,1)\). Los conteos inferiores son \(\lfloor N_j\delta\rfloor-\lfloor N_{j-1}\delta\rfloor\), y dan \((33,11,1,0)\). Los superiores son \(\lceil N_j\delta\rceil-\lceil N_{j-1}\delta\rceil\), con la misma frontera inicial cero, y dan \((34,11,1,0)\). Estos valores dependen del intercepto y del orden de los estratos; no se afirman para una fase inicial arbitraria. Sus totales son cuarenta y cinco y cuarenta y seis. El incremento corresponde a una única unidad de frontera, situada en el primer estrato de esta partición ordenada. Las cifras treinta y cuatro y once aparecen así en un mismo recuento previo a toda carta de unidades; identificarlas con exponentes físicos exige, además, la derivación dimensional correspondiente.

La denominación \emph{cristal aperiódico} se emplea aquí para este orden simbólico, con complejidad, espectro y retorno especificados. La entropía topológica nula no equivale por sí sola a disipación física nula: una realización termodinámica necesita una dinámica energética y una operación de borrado definidas. Del mismo modo, esta lectura esturmiana no sustituye a todos los operadores del TPK ni demuestra por sí sola la periodicidad o aperiodicidad de la salida completa de \(\alpha\). Su función es hacer explícita una estructura de orden que el simple retorno del calendario no describe.


% BEGIN OWNER /Users/ruben/Documents/New project/output/EXTREMA_MEDIA_RAZON_RECIPROCIDAD_ES_EN_20260918/ARTICULO/ES/source/sections/revision_monodromia.tex
% Adición acumulativa. Insertar después de la suspensión esturmiana existente.
% Conservar íntegro el texto y la figura fig:holonomia-esturmiana.
\subsection{Conexión nonádica, holonomía y representación monodrómica}
\label{x_reciprocidad:mono-ampl:conexion}

La conexión regional ya construida se estudia ahora mediante tres objetos:
el transporte del estado completo, la representación de sus vueltas y la
suspensión de sus itinerarios. Las operaciones de APP--TRIT--TPK y el
catálogo inicial conservan las definiciones anteriores. La distinción
necesaria concierne aquí a lo que cada representación retiene de una
misma historia (secciones~\ref{x_reciprocidad:sec:extension} y~\ref{x_reciprocidad:sec:generacion}).

En el índice de prolongación \(k\), una carta del estado conserva
\[
 \widehat x_k=(B_k,C_k,p_k,c_k,\Sigma_k,W_k,g_k,m_k).
\]
Aquí \(B_k\) reúne los tres bloques regionales, \(C_k\) es el cilindro,
\(p_k\) el prefijo, \(c_k\) el registro de acarreo, \(\Sigma_k\) la firma
de frontera y \(W_k\) la información incidencial. La fase es
\(g_k=1+[(k-1)]_9\); \(m_k\) cuenta las vueltas completas. La reducción
\(\beta_k=[m_k]_{12}\) conserva la posición dodecafásica, mientras el
entero \(m_k\) distingue sus sucesivas vueltas.

\begin{definition}[Transporte de una vuelta nonádica]
\label{x_reciprocidad:mono-ampl:def-vuelta}
Sean \(\mathcal R_k\subseteq\widehat X_k\times\widehat X_{k+1}\)
las relaciones de extensión admisible del TPK. La realización de la
holonomía en el nivel \(k\) es la composición relacional
\[
 \Gamma_{9,k}=\mathcal R_{k+8}\circ\cdots\circ\mathcal R_k.
\]
Para \(s\geq1\), su prolongación a \(s\) vueltas es
\[
 \Gamma_{9s,k}=
 \Gamma_{9,k+9(s-1)}\circ\cdots\circ\Gamma_{9,k}.
\]
\end{definition}

La formulación relacional conserva las continuaciones admisibles. Sobre
una historia individualizada, cada flecha lleva al estado efectivamente
prolongado; el bloque visible aislado puede seguir admitiendo más de una
continuación. Esta diferencia es precisamente la función de la memoria y
de las condiciones de frontera en la determinación de la trayectoria.

\begin{proposition}[Iteración del retorno y memoria acumulada]
\label{x_reciprocidad:mono-ampl:vueltas}
En toda composición admisible de \(s\) vueltas,
\[
 g_{k+9s}=g_k,\qquad m_{k+9s}=m_k+s,\qquad
 \beta_{k+9s}=\beta_k+s\pmod{12}.
\]
\end{proposition}
\begin{proof}
La regla de una vuelta conserva la fase e incrementa el contador en una
unidad. Su aplicación al estado ya transportado repite exactamente esas
dos operaciones. La inducción da las dos primeras igualdades; reducir la
segunda módulo \(12\) da la tercera.
\end{proof}

El contador proporciona así un testigo entero de la diferencia entre las
vueltas. Tras \(12\) vueltas la pareja de fases finitas se restaura, pero
el contador aumenta \(12\). La extensión no escindida
\eqref{x_reciprocidad:eq:extension108} expresa la composición del calendario mediante su
cociclo; la trayectoria conserva además la elevación entera de ese dato.
La conservación significa aquí transporte de las contribuciones previas,
no inmovilidad de todas las coordenadas.

\subsubsection{Representación bilateral del índice de retorno}

El desplazamiento \(T\) de \eqref{x_reciprocidad:eq:weyl-relojes} admite una
expresión en coordenadas de fase y vuelta. Numeramos las fases de \(0\) a
\(8\), mediante un desplazamiento del origen respecto de \(g_k\), y
escribimos \(k=9m+r\), con \(0\leq r<9\), en la recta entera
que recubre el ciclo de fases. Sobre
\(\mathcal H_{\rm ind}=\ell^2(\mathbb Z)\), la base
\(|k\rangle=|r,m\rangle\) expresa ese mismo desplazamiento unitario
\[
 T|r,m\rangle=
 \begin{cases}
 |r+1,m\rangle,&r<8,\\
 |0,m+1\rangle,&r=8.
 \end{cases}
\]
La inversa reduce \(r\) en una unidad cuando \(r>0\), y lleva
\(|0,m\rangle\) a \(|8,m-1\rangle\). El paso por el origen de fase
incorpora, por tanto, la unidad exacta de cociente.

\Needspace{8\baselineskip}
\begin{proposition}[Representación de la monodromía del ciclo]
\label{x_reciprocidad:mono-ampl:representacion}
Sea \(S=T^9\). Entonces
\[
 S|r,m\rangle=|r,m+1\rangle,\qquad
 \mathbb Z\longrightarrow\mathcal U(\mathcal H_{\rm ind}),
 \quad s\longmapsto S^s
\]
es una representación unitaria fiel del grupo de vueltas del ciclo.
\end{proposition}
\begin{proof}
Nueve desplazamientos conservan el resto de la división por \(9\) e
incrementan el cociente. Las potencias satisfacen
\(S^{s+t}=S^sS^t\), y \(S^s|r,m\rangle=|r,m+s\rangle\) distingue
todo entero \(s\). La permutación de una base ortonormal es unitaria.
\end{proof}

Esta representación corresponde al recubrimiento del ciclo, cuyo grupo
de vueltas es \(\mathbb Z\), y representa la fase con su contador. El
estado del TPK contiene además las coordenadas de trayectoria ya descritas.
La extensión bilateral admite índices negativos como coordenadas del
recubrimiento; una trayectoria iniciada en una condición inicial utiliza
su semieje admisible. De este modo se dispone de una inversa operatoria
del índice sin convertir cada relación de extensión en una biyección del
estado completo.

Con \(D_M|k\rangle=\lfloor k/9\rfloor|k\rangle\), la identidad
\([D_M,S]=S\) expresa el incremento de memoria en el dominio denso de
vectores de soporte finito. Si \(\mathcal I|k\rangle=|-k\rangle\) y
\(P_0\) proyecta sobre los índices divisibles por \(9\), se obtiene
además
\begin{equation}
 \mathcal I T\mathcal I=T^{-1},\qquad
 \mathcal I D_M\mathcal I=-D_M-(I-P_0).
 \label{x_reciprocidad:mono-ampl:inversion-contador}
\end{equation}
En efecto, para \(k=9m+r\),
\(\lfloor-k/9\rfloor=-m\) si \(r=0\) y vale \(-m-1\) si \(r>0\).
La inversión conserva así el término de frontera de la división euclídea.

\subsection{Reloj de bloques y reloj de capacidad}
\label{x_reciprocidad:mono-ampl:relojes}

El índice de una aplicación de refinamiento y el número de posiciones
decimales disponibles tienen leyes relacionadas, pero diferentes. Para
expresarlas reservamos \(n\) al índice de bloques ternarios y definimos
\[
 \rho=\log_{1000}729=\log_{10}9,\qquad
 \delta=1-\rho,\qquad
 \mathcal C_n=\lfloor(n+1)\rho\rfloor,\quad n\geq0.
\]
La letra \(\mathcal C_n\) designa aquí la capacidad, no el registro
dodecafásico \(K\). En el origen de fase canónico, la palabra mecánica
anterior proporciona
\[
 Z_n=\sum_{j=0}^{n}z_j(0)=\lfloor(n+1)\delta\rfloor.
\]

\begin{proposition}[Balance exacto entre capacidad y vacancias]
\label{x_reciprocidad:mono-ampl:balance-capacidad}
Para \(n\geq0\) y, en la segunda identidad, \(n\geq1\),
\[
 \mathcal C_n=n-Z_n,\qquad
 \mathcal C_n-\mathcal C_{n-1}=1-z_n(0).
\]
\end{proposition}
\begin{proof}
La irracionalidad de \(\delta\) implica
\(\lceil(n+1)\delta\rceil=\lfloor(n+1)\delta\rfloor+1\).
Entonces
\[
 \lfloor(n+1)(1-\delta)\rfloor
 =n+1-\lceil(n+1)\delta\rceil=n-Z_n.
\]
Restar dos identidades consecutivas concluye la prueba.
\end{proof}

Un evento \(z_n=1\) añade un bloque a la historia manteniendo la capacidad
\(\mathcal C_n\); un evento \(z_n=0\) incrementa ambos contadores. En
cualquier intervalo, los incrementos de capacidad y las vacancias suman
exactamente su longitud. Reduciendo módulo \(q\), para \(q=9\) o
\(q=108\), resulta
\[
 [\mathcal C_n]_q=[n]_q-[Z_n]_q.
\]
La fase uniforme de bloques y la fase de capacidad se relacionan mediante
la memoria integrada de vacancias. En particular, una retención de
capacidad y una vuelta completa de la conexión tienen mecanismos
distintos, aunque ambas puedan conservar una observación de fase.

La inversión monótona del reloj da un tiempo de primera llegada:
\[
 T_{\rm cap}(a)=\min\{n:\mathcal C_n=a\}
 =\left\lceil\frac a\rho\right\rceil-1,
 \qquad a\geq1.
\]
Con \(\sigma=\rho^{-1}-1\), el tiempo empleado en ganar \(b\) unidades
de capacidad es
\[
 T_{\rm cap}(a+b)-T_{\rm cap}(a)
 =b+\lceil(a+b)\sigma\rceil-\lceil a\sigma\rceil.
\]
La diferencia de techos toma uno de los dos enteros adyacentes a
\(b\sigma\). Así se obtienen \(9\) o \(10\) bloques para ganar
\(9\) unidades de capacidad, y \(113\) o \(114\) para ganar \(108\).
Son tiempos del reloj inverso, medidos desde una primera llegada; el
contador de vueltas se interpreta siempre en el índice al que pertenece.

\subsection{Suspensión simbólica y cristal temporal aperiódico}
\label{x_reciprocidad:mono-ampl:cristal}

La codificación por rotación reúne periodicidad de fase, recurrencia local
y aperiodicidad global. Sea \(\mathcal X_\delta\subseteq\{0,1\}^{\mathbb Z}\)
la clausura, en la topología producto, de las traslaciones de la palabra
bilateral \(z_n(\theta)\). Denotemos por \(\sigma_{\rm sh}\) el
desplazamiento de sus sucesiones. Conservando una fase nonádica, el paso es
\[
 F(r,z)=([r+1]_9,\sigma_{\rm sh}z),\qquad
 (r,z)\in C_9\times\mathcal X_\delta.
\]

\begin{definition}[Cristal temporal aperiódico simbólico]
\label{x_reciprocidad:mono-ampl:def-cristal}
La suspensión de techo unitario del sistema anterior es el espacio
\[
 \mathcal S_\delta=
 \bigl((C_9\times\mathcal X_\delta)\times\mathbb R\bigr)/\sim,
 \qquad (y,t+1)\sim(Fy,t),
\]
con flujo \(\Phi_s[y,t]=[y,t+s]\). Se denomina \emph{cristal temporal
aperiódico simbólico} a este modelo de orden temporal, junto con su
codificación de vacancias y la representación del contador de vueltas.
\end{definition}

El adjetivo temporal designa la acción del flujo y de su sección de retorno;
la aperiodicidad caracteriza sus itinerarios. La frecuencia de unos de cada
sucesión es \(\delta\): la cota uniforme de discrepancia menor que uno,
obtenida por telescopía en cualquier ventana, persiste al tomar la clausura.
Una sucesión periódica tendría frecuencia racional. Por tanto,
\(\mathcal X_\delta\) no contiene órbitas periódicas. La suspensión
tampoco tiene órbitas cerradas: un retorno del flujo exigiría un tiempo
entero y una periodicidad del itinerario.

Las palabras finitas sí reaparecen. Los itinerarios de longitud \(m\)
proceden de una partición de la circunferencia por \(m+1\) puntos; cada
intervalo es visitado recurrentemente por la rotación irracional. Incluso
al restringir los tiempos a una fase fija, el paso \(9\delta\) sigue
siendo irracional. Ésta es la razón de que cada palabra aparezca en las
\(9\) fases y de las complejidades ya establecidas
\[
 p(m)=m+1,\qquad p_9(m)=9(m+1),\qquad p_3(m)=3(m+1).
\]
La recurrencia de configuraciones locales coexiste con la ausencia de un
período global. Su entropía topológica es nula porque estas complejidades
crecen linealmente.

La realización por los observables \(V_9,W_\delta\) y el desplazamiento
\(T\) expresa el mismo orden mediante las covariancias
\eqref{x_reciprocidad:eq:weyl-relojes}. Las frecuencias del reloj uniforme pertenecen a
\(\frac19\mathbb Z+\delta\mathbb Z\), módulo \(1\). El reloj de
capacidad conserva adicionalmente la relación
\([\mathcal C_n]_9=[n]_9-[Z_n]_9\); no se sustituye por la fase uniforme
al interpretar los retornos. El término \emph{cristal temporal} queda así
definido mediante un sistema dinámico matemático. Su realización física
requiere especificar una dinámica energética, sus observables temporales
y su estabilidad, datos distintos de la construcción simbólica presente.

\subsection{Simetría especular de las regiones y doble realización}
\label{x_reciprocidad:mono-ampl:regiones}

El intercambio especular actúa en el bloque regional
\(B=(u^\pi,u^e,u^\varphi)\in(\mathbb F_3^6)^3\) por
\[
 \mathfrak cB=(u^e,u^\pi,u^\varphi).
\]
Su eje fijo y su agregado son, respectivamente, \(u^\varphi\) y
\(s=u^\pi+u^e\). Esta fijación se refiere a la involución en una fibra;
ambas coordenadas pueden evolucionar durante el transporte nonádico.
Con \(q=u^\pi+u^e+u^\varphi\),
\(a=u^\pi+u^e-u^\varphi\) y \(d=u^\pi-u^e\), se recupera
\[
 u^\varphi=2(q-a),\qquad s=q-u^\varphi,\qquad
 u^\pi=2(s+d),\qquad u^e=2(s-d).
\]
Las igualdades son componente a componente en \(\mathbb F_3\), donde
\(2^{-1}=2\). El espejo conserva \(q,a,s\) e invierte \(d\); esta
coordenada antisimétrica individualiza el reparto que el agregado conserva
sin distinguir. La inversión \(\mathcal I\) del índice temporal y la
involución \(\mathfrak c\) de canales actúan, por consiguiente, sobre
datos diferentes.

El retorno de fase posee un testigo explícito:
\[
 \begin{aligned}
 B_1&=(012222\mid121221\mid112202),\\
 B_{10}&=(101222\mid112221\mid112002).
 \end{aligned}
\]
Las dos posiciones tienen fase \(1\); el bloque, el eje y la memoria han
cambiado. Asimismo, la sincronización de los censos en las fases \(8\) y
\(9\) anula sus diferencias relativas, mientras la componente diagonal
pasa de \((59,59,59)\) a \((43,43,43)\). La igualdad de un observable
relativo no agota el estado que lo determina.

La doble realización conserva esa arquitectura. En la sección
\ref{x_reciprocidad:exc:doble-lectura}, el mismo prefijo determina por una parte cilindros
encajados y por otra una sucesión de palabras codificadas con sus marcos.
El primer mapa pasa al límite por contracción de diámetros; el segundo
satisface \(r_{n,m}Q_n=Q_mp_{n,m}\) y define \(Q_\infty\).
La figura \ref{x_reciprocidad:mono-ampl:fig-regional} anticipa esas dos aplicaciones.
El recorrido incidencial posterior construye códigos, retículos y la
realización de Moonshine en sus dominios propios; el grupo de automorfismos
actúa sobre esa realización, no sobre las cifras por una identificación
implícita.


% BEGIN OWNER /Users/ruben/Documents/New project/output/EXTREMA_MEDIA_RAZON_RECIPROCIDAD_ES_EN_20260918/ARTICULO/ES/source/figures/monodromia_regional_ampliacion.tex
% Figura acumulativa; destino figures/monodromia_regional_ampliacion.tex.
\begin{figure}[p]
\centering
\begin{tikzpicture}[x=1cm,y=1cm,>=Latex,font=\small,
  flow/.style={->,line width=.55pt},
  ann/.style={align=center,font=\scriptsize}]
\node[font=\bfseries] at (6.9,9.2)
  {Conexión nonádica y coordenadas regionales};
\begin{scope}[shift={(3.4,6.8)}]
 \foreach \k in {1,...,9}{
  \pgfmathsetmacro{\ang}{90-40*(\k-1)}
  \coordinate (p\k) at (\ang:1.52);
  \fill (p\k) circle (1.4pt);
  \node[fill=white,inner sep=1.1pt,font=\scriptsize]
    at (\ang:1.83) {\(\k\)};
 }
 \foreach \k in {1,...,8}{
  \pgfmathtruncatemacro{\next}{\k+1}
  \draw[flow,shorten >=3pt,shorten <=3pt] (p\k)--(p\next);
 }
 \draw[flow,shorten >=3pt,shorten <=3pt] (p9)--(p1);
 \node[ann] at (0,.1) {\(g_k\in C_9\)\\\(\Gamma_{9,k}\)};
 \node[ann] at (0,-2.23) {Una vuelta: \(g\mapsto g\), \(m\mapsto m+1\)};
\end{scope}
\node[ann,text width=5.25cm] at (10.1,7.75)
 {\(B_k=(u_k^\pi,u_k^e,u_k^\varphi)\)\\
 bloque regional en una fibra};
\node at (8.55,6.5) {\(u^\pi\)};
\node at (11.65,6.5) {\(u^e\)};
\draw[<->,line width=.65pt] (8.95,6.5)--node[above]{\(\mathfrak c\)}(11.25,6.5);
\draw[densely dashed,gray] (10.1,5.72)--(10.1,7.23);
\node[fill=white,inner sep=2pt] at (10.1,5.6){\(u^\varphi\)};
\node[ann] at (10.1,4.6)
 {Eje \(u^\varphi\) y agregado \(u^\pi+u^e\)\\
 invariantes bajo \(\mathfrak c\)};
\draw[flow] (5.45,6.8)--(7.25,6.8);
\draw[gray!60,line width=.3pt] (.55,4.02)--(13.25,4.02);
\node[ann,text width=12cm] (hist) at (6.9,3.47)
 {Historia compatible: prefijos, cilindros, orientación, cocientes,
 marcos de incidencia y memoria};
\coordinate (fork) at (6.9,2.85);
\draw[flow] (hist.south)--(fork);
\draw[flow] (fork)--(3.6,2.2);
\draw[flow] (fork)--(10.3,2.2);
\node[ann,text width=5.6cm] at (3.6,1.65)
 {Realización arquimediana\\
 \(I_{n+1}\subseteq I_n,\quad |I_n|=729^{-n}\)\\
 límite de los prefijos regionales};
\node[ann,text width=5.6cm] at (10.3,1.65)
 {Realización incidencial\\
 \(r_{n,m}Q_n=Q_mp_{n,m}\)\\
 palabras codificadas y marcos compatibles};
\node[ann,text width=5.6cm] at (3.6,.42)
 {\(\pi,\varphi,e\); determinación dodecafásica de \(\alpha\)};
\node[ann,text width=5.6cm] at (10.3,.42)
 {Witt--Golay--Leech--\(V^\natural\)\\
 automorfismos de la realización excepcional};
\draw[flow] (3.6,1.03)--(3.6,.78);
\draw[flow] (10.3,1.03)--(10.3,.78);
\end{tikzpicture}
\caption{Nueve fases de una conexión y dos realizaciones de la historia
compatible. El ciclo representa la fase; cada vuelta incrementa el contador.
El diagrama regional representa la involución \(\mathfrak c\), que
intercambia clausura y propagación y fija el eje de autoescala en una fibra.
Esta invariancia no impone que el eje sea constante durante el transporte.
Las dos aplicaciones inferiores utilizan el mismo prefijo: una determina
su límite numérico y la otra conserva palabras y marcos de incidencia. Los
mapas de esta segunda realización se especifican en la sección
\ref{x_reciprocidad:exc:doble-lectura}; el recorrido excepcional se desarrolla después.}
\label{x_reciprocidad:mono-ampl:fig-regional}
\end{figure}

% END OWNER



% BEGIN OWNER /Users/ruben/Documents/New project/output/EXTREMA_MEDIA_RAZON_RECIPROCIDAD_ES_EN_20260918/ARTICULO/ES/source/sections/recta_nonadica_recuperada.tex
\clearpage
\subsubsection{Representación rectilínea de las fases y del retorno}
\label{x_reciprocidad:mono-recta:desarrollo}

La representación sobre el segmento \([0,10]\) conserva el origen
posicional de APP y ordena sus nueve marcas interiores. En este diagrama,
los enteros \(1,\ldots,9\) son representantes de fase. La posición
horizontal indica el orden de las cartas; las etiquetas superiores
especifican operaciones regionales. La figura
\ref{x_reciprocidad:mono-recta:figura} reúne así la recta de partida, la fijación del
eje de autoescala, la distribución especular de los dos canales y el
retorno con memoria (sección~\ref{x_reciprocidad:mono-ampl:conexion}).


% BEGIN OWNER /Users/ruben/Documents/New project/output/EXTREMA_MEDIA_RAZON_RECIPROCIDAD_ES_EN_20260918/ARTICULO/ES/source/figures/recta_nonadica_0_10.tex
% Recuperación de fig_nonadic_route_mini.tex de la síntesis académica.
% Los rótulos designan operaciones regionales; la suma se efectúa en F_3^6.
\begin{figure}[H]
\centering
\begin{tikzpicture}[x=1.12cm,y=1cm,font=\small,>=Latex,
  guide/.style={line width=.45pt,black!55}]
 \draw[line width=.8pt] (0,0)--(10,0);
 \foreach \k in {0,...,10}{
  \draw[guide] (\k,-.09)--(\k,.09);
  \node[below=2.5mm] at (\k,0) {$\k$};
 }
 \foreach \k in {1,...,9}{\fill (\k,0) circle (1.7pt);}
 \node[anchor=west,font=\footnotesize] at (0,3.05)
   {Segmento generador $[0,10]$ y nueve representantes de fase};
 \draw[guide] (5,.12)--(5,1.7);
 \node[above,align=center,font=\footnotesize] at (5,1.7)
   {fijación del eje\\$u^\varphi$};
 \draw[decorate,decoration={brace,amplitude=4pt},guide]
   (6,.65)--(8,.65);
 \node[above=5pt,align=center,font=\footnotesize] at (7,.65)
   {agregado regional\\$u^e+u^\pi$};
 \draw[guide] (9,.12)--(9,1.7);
 \node[above,align=center,font=\footnotesize] at (9,1.7)
   {involución especular\\$\pi\leftrightarrow e$};
 \draw[->,line width=.7pt] (9,-.7)
   .. controls (9,-1.55) and (1,-1.55) .. (1,-.7);
 \node[fill=white,inner sep=3pt,font=\footnotesize] at (5,-1.23)
   {retorno $9\to1$;\quad $m\mapsto m+1$};
 \node[font=\footnotesize] at (5,-1.95)
   {$w_6\to w_{12}\to w_{18}\to w_{24}\to w_{30}\to R_{36}\to G_9$};
\end{tikzpicture}
\caption{Representación rectilínea de la conexión nonádica. La marca
\(5\) sitúa la primera saturación fuerte y el eje regional de
autoescala; la llave \(6\)--\(8\) identifica el agregado ternario
fijado en cada fibra; la fase \(9\) selecciona la orientación de
clausura. El retorno conserva la fase e incrementa la memoria. Las
abscisas numeran posiciones operativas, mientras las etiquetas
\(u^\pi,u^e,u^\varphi\) designan bloques de \(\mathbb F_3^6\).}
\label{x_reciprocidad:mono-recta:figura}
\end{figure}

% END OWNER


La fijación del eje se formula en la fibra de una frontera dada. Si
\(q=u^\pi+u^e+u^\varphi\) y
\(a=u^\pi+u^e-u^\varphi\), entonces
\[
 u^\varphi=2(q-a),\qquad s=u^\pi+u^e=q-u^\varphi
 \quad\text{en }\mathbb F_3^6.
\]
Los datos \((q,a)\) determinan el eje y el agregado; el reparto ordenado
de \(s\) conserva la libertad especular resuelta por la prolongación.
La quinta fase posee una solución local y una extensión, con dato
de frontera \(c=111111\) y firma residual \(\rho_W=000\). Éste es el
sentido de su primera saturación fuerte. El símbolo \(\varphi\)
situado sobre la marca \(5\) identifica el eje regional que interviene
en esa operación.

Las fases intermedias conservan esta descomposición dentro de cada fibra,
mientras sus datos de frontera evolucionan. El censo completo de
multiplicidades locales y extensiones de la rama distinguida es
\[
\begin{array}{c|rrrrrrrrr}
\text{fase}&1&2&3&4&5&6&7&8&9\\\hline
N_{\rm loc}&3&2&5&4&1&1&3&3&2\\
N_{\rm ext}&2&5&4&1&1&3&3&2&1
\end{array}
\]
La fase \(6\) combina unicidad local y tres prolongaciones; la fase
\(7\) resuelve una fibra especular triple; la fase \(8\) sincroniza
los tres censos regionales. La fase \(9\) conserva dos distribuciones
locales con el mismo eje y el mismo agregado, de las cuales una se
prolonga. Estos recuentos expresan funciones diferentes dentro de una
misma conexión, y justifican las marcas de la recta.

En particular, para las fases \(6,7,8\) las sumas regionales son
respectivamente \(012100\), \(101001\) y \(112000\). La conservación
representada por la llave es una invariancia respecto del reparto
especular en cada fibra. La evaluación real posterior aplica los
lectores de los cilindros compatibles a las historias completas; la
suma interna representada en la figura pertenece al dominio ternario.

Finalmente, el arco \(9\to1\) representa el paso de una vuelta a la
siguiente. Las fórmulas de la sección
\ref{x_reciprocidad:mono-ampl:conexion} conservan el resto de fase e incrementan el
contador de vueltas. El diagrama rectilíneo y el diagrama cíclico de la
figura \ref{x_reciprocidad:mono-ampl:fig-regional} expresan, por tanto, dos aspectos
complementarios de la misma holonomía: orden de las posiciones y
composición del retorno sobre el estado con memoria.

% END OWNER


\subsection{Transporte temporal y composición exponencial trítica}
\label{x_reciprocidad:mono-ampl:euler}

Las relaciones de Euler de la sección \ref{x_reciprocidad:euler-trit-articulacion}
describen las realizaciones locales de los tres regímenes. Para una firma
fija \(\tau\), el generador satisface \(J_\tau^2=-\tau I\); la
conjugación algebraica lleva \(J_\tau\) a \(-J_\tau\). En consecuencia,
\[
 \exp(sJ_\tau)\exp(tJ_\tau)=\exp((s+t)J_\tau),\qquad
 \overline{\exp(tJ_\tau)}=\exp(-tJ_\tau),
\]
y el producto de una evolución con su conjugada es la unidad. En el tipo
elíptico se realiza mediante \(\cos^2t+\sin^2t=1\); en el hiperbólico,
mediante \(\cosh^2t-\sinh^2t=1\); en el parabólico,
\((I+tN)(I-tN)=I\), pues \(N^2=0\). La misma ley de conservación
admite tres expresiones cuadráticas.

La realización elíptica cierra una vuelta; la nilpotente conserva un
desplazamiento lineal; la hiperbólica separa expansión y contracción con
determinante unitario. En los parámetros regionales ya generados,
\[
 \exp(\pi J_{\mathrm e})=-I,\qquad
 \exp(J_{\mathrm p})=I+J_{\mathrm p},\qquad
 \exp(2\log\varphi\,J_{\mathrm h})=F_{\rm av}^{\,2}.
\]
La evaluación \(\exp(I)=eI\) reconoce la unidad de la ley exponencial
común. El número \(e\) no se asigna exclusivamente al régimen parabólico.

El transporte de una realización local por un isomorfismo \(U\) conserva
estas relaciones: si \(J'=UJU^{-1}\), la serie exponencial da
\(U\exp(tJ)=\exp(tJ')U\). La identidad transporta el régimen y su
norma. Una transición entre tipos cuadráticos diferentes pertenece, en
cambio, a la selección de hojas y regímenes del TPK; no es una conjugación
que cambie el valor de \(J^2\). El producto temporal conserva por ello el
orden de las transiciones y la información de sus extremos.

Esta articulación distingue la clausura de una representación local y la
holonomía de una historia. Puede cumplirse \(\exp(2\pi J_{\mathrm e})=I\)
al mismo tiempo que \(S|r,m\rangle=|r,m+1\rangle\). La primera igualdad
restaura la orientación del régimen; la segunda registra otra vuelta. La
estructura discreta conserva ambas: el retorno local y la procedencia
acumulada del estado que lo realiza.

% END OWNER



% BEGIN OWNER /Users/ruben/Documents/New project/output/EXTREMA_MEDIA_RAZON_RECIPROCIDAD_ES_EN_20260918/ARTICULO/ES/source/sections/primos_retornos.tex
% Recuperación del producto nativo y su selector, no incorporación de RH.
\subsection{Irreducibilidad multiplicativa y períodos de retorno}
\label{x_reciprocidad:primos-retornos}

La distinción entre estructura, evaluación y fase aparece también en el
sector multiplicativo. El producto local de APP contiene tanto el residuo
como el cociente; su iteración permite construir formas normales y
factorizaciones antes de examinar sus períodos decimales. Este desarrollo
sitúa la expresión \emph{modo primo} en un dominio preciso y relaciona la
irreducibilidad con las operaciones ya definidas, no con una coincidencia
aislada de períodos (sección~\ref{x_reciprocidad:primos-retornos}).

Sea
\[
 \mathcal W_9=\{\varnothing\}\sqcup
 \bigsqcup_{\ell\geq1}\{1,\ldots,9\}^{\ell}.
\]
Las secuencias se ordenan desde el dígito menos significativo. La evaluación
\(\operatorname{ev}_9(u)=\sum_j u_j9^j\), con
\(\operatorname{ev}_9(\varnothing)=0\), es biyectiva sobre los enteros no
negativos: para un entero positivo, la división de \(n-1\) por \(9\)
determina de manera única el primer dígito y la continuación. Esta es la
numeración nonádica biyectiva; conserva el representante \(9\) de la
frontera.

La suma \(\boxplus_\#\) se obtiene aplicando la hoja aditiva de APP a los
dígitos de cada posición y propagando sus cocientes. Si sólo hay un dígito,
éste constituye el residuo provisional; un cociente entrante no nulo se
normaliza mediante una segunda operación aditiva. La suma de los
cocientes se transporta a la posición siguiente. Para multiplicar una
secuencia por un dígito \(d\), la celda multiplicativa produce
\(u_jd=9q_j+r_j\), y el cociente entrante \(c_j\) se incorpora mediante
\[
 (w_j,c_{j+1})=
 \begin{cases}
 (r_j,q_j),&c_j=0,\\
 \bigl(R^+(r_j,c_j),q_j+q^+(r_j,c_j)\bigr),&c_j>0.
 \end{cases}
\]
Se comienza con \(c_0=0\) y se incorpora el cociente final cuando es no
nulo. Durante la suma, \(c_j\leq2\); durante el producto por un dígito,
\(c_j\leq9\). Las operaciones locales reciben así argumentos dentro
de su dominio. Denotemos por \(d\odot_\#u\) este segundo procedimiento,
con \(d\odot_\#\varnothing=\varnothing\).

Para \(v=(v_0,\ldots,v_{m-1})\), el producto completo se construye por
la recurrencia
\[
 a_m=\varnothing,\qquad
 a_j=(9\odot_\#a_{j+1})\boxplus_\#(v_j\odot_\#u),
 \qquad u\boxtimes_\#v=a_0.
\]
La regla compone las dos hojas locales. Su evaluación verifica
\begin{equation}
 \operatorname{ev}_9(u\boxtimes_\#v)
   =\operatorname{ev}_9(u)\operatorname{ev}_9(v).
 \label{x_reciprocidad:eq:primos-entrelazamiento}
\end{equation}
En efecto, cada transición conserva la suma parcial más el cociente
pendiente ponderado por la potencia de \(9\). La recurrencia da entonces
\(\operatorname{ev}_9(a_j)=9\operatorname{ev}_9(a_{j+1})+
v_j\operatorname{ev}_9(u)\), cuya iteración prueba
\eqref{x_reciprocidad:eq:primos-entrelazamiento}. La unicidad de la forma normal transporta
asociatividad y conmutatividad al producto construido; su unidad es \((1)\).

\begin{definition}[Coproducto multiplicativo reducido]
Sobre el espacio vectorial de soporte finito generado por las formas
normales no vacías se define
\[
 \widetilde\Delta_\# e_w
 =\sum_{\substack{u\boxtimes_\#v=w\\u,v\ne(1)}}e_u\otimes e_v.
\]
Una forma normal no unitaria es irreducible cuando su imagen por
\(\widetilde\Delta_\#\) es nula.
\end{definition}

Las fibras de factorización son finitas por
\eqref{x_reciprocidad:eq:primos-entrelazamiento}. En la completación
\(\mathcal H_\#=\ell^2(\mathcal W_9\setminus\{\varnothing\})\),
las imágenes de formas distintas tienen soportes disjuntos. Si
\(d_\#(w)\) cuenta las factorizaciones ordenadas no unitarias, el dominio
de la clausura es
\[
 \operatorname{Dom}(\overline{\widetilde\Delta_\#})
 =\left\{x\in\mathcal H_\#:
       \sum_w d_\#(w)|x_w|^2<\infty\right\}.
\]
En efecto, la norma cuadrada de la imagen es esa suma ponderada. Por ello,
el operador
\[
 A_\#=(\overline{\widetilde\Delta_\#})^*
                  \overline{\widetilde\Delta_\#}
\]
es diagonal, positivo y autoadjunto: su valor propio en \(e_w\) es el
número de factorizaciones ordenadas no unitarias de \(w\). La proyección
\[
 P_{\rm irr}=(I-|e_{(1)}\rangle\langle e_{(1)}|)
                    \mathbf1_{\{0\}}(A_\#)
\]
selecciona exactamente las formas irreducibles. Mediante la evaluación
multiplicativa, éstas corresponden a los primos ordinarios. La selección
depende de la estructura de factorización, y mantiene separado el estado
unitario.

Para un valor \(n\geq2\) coprimo con \(10\), la división decimal induce la
permutación de residuos \(r\mapsto10r\pmod n\). El retorno del residuo
\(1\) tiene período
\[
 \tau_n=\operatorname{ord}_n(10).
\]
Si \(n=\prod_p p^{a_p}\), el teorema chino de los restos da
\[
 \tau_n=\operatorname{mcm}_{p^{a_p}\parallel n}
                    \operatorname{ord}_{p^{a_p}}(10).
\]
El retorno compuesto sincroniza los retornos de las potencias primas.
Para un primo \(p\notin\{2,3,5\}\), se tiene
\[
 \tau_p\mid p-1,\qquad
 M_p=\frac{10^{\tau_p}-1}{p}\in\mathbb N,
 \qquad M_p\equiv0\pmod9.
\]
La última congruencia expresa el cierre nonádico de la repetición decimal:
\(10^{\tau_p}-1\) es múltiplo de \(9\) y \(p\) es invertible módulo
\(9\). La misma congruencia vale para cualquier \(n\) coprimo con
\(90\). En particular, \(7,13,77,91,143\) tienen período decimal \(6\),
aunque sólo los dos primeros son irreducibles. El coproducto y el período
registran, por tanto, propiedades diferentes de un mismo valor construido.

Esta diferenciación precisa la relación con el refinamiento aperiódico.
La secuencia de vacancias determina los incrementos del índice de
capacidad; la factorización determina los modos multiplicativos; la
división decimal determina sus tiempos de retorno. El representante
nonádico conserva el cierre de fase, mientras que el cociente y la
factorización conservan información que ese cierre no distingue. El
enlace con las construcciones espectrales del tratado reside en esos
operadores y sus entrelazamientos. La demostración de sus resultados
analíticos posteriores mantiene su desarrollo propio; no se sustituye por
la periodicidad de una observación residual.

\subsubsection{Discrepancia de vacancias y lecturas de Dirichlet}

El vínculo admite también una expresión analítica exacta. Se conserva la
misma codificación \(z_n=z_n(0)\), para \(n\geq1\), y se consideran
dos observables posteriores: la serie sobre todos los índices y el
producto sobre los índices seleccionados por irreducibilidad. Para
\(\operatorname{Re}s>1\),
\[
 Z_{\rm vac}^{+}(s)=\sum_{n\geq1}\frac{z_n}{n^s},\qquad
 Z_{\rm vac}^{\times}(s)=\prod_p(1-p^{-s})^{-z_p}.
\]
Ambos convergen absolutamente en ese semiplano. Su factor común es la
densidad \(\delta\), pero sus discrepancias son diferentes.

\Needspace{21\baselineskip}
\begin{proposition}[Separación de la densidad de vacancias]
Se tiene
\[
 Z_{\rm vac}^{+}(s)=\delta\zeta(s)+H_{\rm vac}(s),
 \qquad H_{\rm vac}\in\mathcal O(\operatorname{Re}s>0).
\]
En \(\operatorname{Re}s>1\), con logaritmos definidos por las series
absolutamente convergentes de Euler,
\[
 Z_{\rm vac}^{\times}(s)=\zeta(s)^\delta G_{\rm vac}(s),\qquad
 \log G_{\rm vac}(s)=
 \sum_p(z_p-\delta)\log(1-p^{-s})^{-1}.
\]
\end{proposition}
\begin{proof}
Las sumas parciales centradas satisfacen
\[
 A_N:=\sum_{n=1}^N(z_n-\delta)
 =\lfloor(N+1)\delta\rfloor-N\delta,\qquad |A_N|<1.
\]
La sumación de Abel representa el término centrado por
\(s\int_1^\infty A_{\lfloor x\rfloor}x^{-s-1}\,dx\). La integral
converge localmente de manera uniforme en \(\operatorname{Re}s>0\),
lo que prueba la holomorfía. La segunda identidad resulta de separar
\(z_p=\delta+(z_p-\delta)\) en el logaritmo del producto, dentro de
su dominio de convergencia absoluta.
\end{proof}

La descomposición hace visible qué añade la lectura sobre primos: examina
la misma sucesión de vacancias sobre un subconjunto determinado por la
factorización. El control de la discrepancia sobre todos los enteros da
la primera prolongación holomorfa; el segundo observable conserva una
discrepancia propia sobre primos. Son dos aplicaciones de la misma
estructura discreta, con dominios analíticos expresos. Esta articulación
explica su pertinencia en el núcleo aritmético sin trasladar a este
artículo la demostración de los resultados terminales del tratado.

% END OWNER

